\documentclass[12pt]{article}
\usepackage[utf8]{inputenc}
\usepackage{float}
\usepackage{amsmath}

\usepackage{tikz}
\usepackage[hmargin=3cm,vmargin=6.0cm]{geometry}
\topmargin=-2cm
\addtolength{\textheight}{6.5cm}
\addtolength{\textwidth}{2.0cm}
\setlength{\oddsidemargin}{0.0cm}
\setlength{\evensidemargin}{0.0cm}
\usepackage{indentfirst}
\usepackage{amsfonts}

\begin{document}

\section*{Student Information}

Name : Alperen OVAK \\

ID : 2580801\\


\section*{Answer 1}
\subsection*{a)} 

We need to use the Normal approximation to determine the sample size \( n \) for the Monte Carlo simulation such that with 0.99 probability, the estimate differs from the true value by no more than 0.02.

According to our book chapter 5.3, size of a Monte Carlo study is given by:

\[ N \geq p(1-p){\left(\frac{z_{\alpha/2}}{\varepsilon}\right) }^2 \]

Thus, the margin of error \( \varepsilon  \) for a Monte Carlo simulation is:

\[ \varepsilon  = z_{\alpha/2} \sqrt{\frac{p(1-p)}{N}} \]

Where:
\begin{itemize}
    \item \( z_{\alpha/2} \) is the critical value from the standard normal distribution for a 99\% confidence level. For 99\% confidence, \( \alpha = 0.01 \), so \( \alpha/2 = 0.005 \), and \( z_{\alpha/2} \approx 2.576 \).
    \item \( p \) is the estimated probability (unknown, but we can use \( p = 0.5 \) for the most conservative estimate).
    \item \( N \) is the sample size.
\end{itemize}


We want the margin of error \( \varepsilon  \) to be less than 0.02:

\[ 0.02 = 2.576 \sqrt{\frac{0.5(1-0.5)}{N}} \]

Solving for \( N \):

\[ 0.02 = 2.576 \sqrt{\frac{0.25}{N}} \]

\[ 0.02 = 2.576 \cdot \frac{0.5}{\sqrt{N}} \]

\[ 0.02 = 1.288 / \sqrt{N} \]

\[ \sqrt{N} = 1.288 / 0.02 \]

\[ \sqrt{N} = 64.4 \]

\[ N \approx 64.4^2 \]

\[ N \approx 4147.36 \]

Round up:

\[ N = 4148 \]

Thus, the size of the Monte Carlo simulation should be at least 4148.

\subsection*{b)} 

\begin{enumerate}
    \item \textbf{Expected value for the weight of an automobile:}
    
    The weight of each automobile is Gamma distributed with parameters \(\alpha = 120\) and \(\lambda = 0.1\).
    
    The expected value \( E[X] \) for a Gamma distribution is given by:
    
    \[
    E[X] = \frac{\alpha}{\lambda}
    \]
    
    For the automobile:
    
    \[
    E[\text{Weight of an automobile}] = \frac{120}{0.1} = 1200 \, \text{kg}
    \]
  
    \item \textbf{Expected value for the weight of a truck:}
    
    The weight of each truck is Gamma distributed with parameters \(\alpha = 14\) and \(\lambda = 0.001\).
    
    For the truck:
    
    \[
    E[\text{Weight of a truck}] = \frac{14}{0.001} = 14000 \, \text{kg}
    \]
  
    \item \textbf{Expected value for the total weights of all automobiles:}
    
    The number of automobiles passing over the bridge per day is a Poisson random variable with \(\lambda_A = 60\).
    
    Since the weight of each automobile is independent and identically distributed, the total weight of all automobiles is:
    
    \[
    E[\text{Total weight of all automobiles}] = E[N_A] \cdot E[\text{Weight of an automobile}]
    \]
    
    Where \( N_A \) is the number of automobiles:
    
    \[
    E[\text{Total weight of all automobiles}] = 60 \cdot 1200 = 72000 \, \text{kg}
    \]
  
    \item \textbf{Expected value for the total weights of all trucks:}
    
    The number of trucks passing over the bridge per day is a Poisson random variable with \(\lambda_T = 12\).
    
    Since the weight of each truck is independent and identically distributed, the total weight of all trucks is:
    
    \[
    E[\text{Total weight of all trucks}] = E[N_T] \cdot E[\text{Weight of a truck}]
    \]
    
    Where \( N_T \) is the number of trucks:
    
    \[
    E[\text{Total weight of all trucks}] = 12 \cdot 14000 = 168000 \, \text{kg}
    \]
  \end{enumerate}
  

\section*{Answer 2}

\subsection*{Simulation Size}
We use \( N = 4148 \) as determined from part (a) to ensure that our simulation estimate is accurate within the desired margin of error.

\subsection*{Generating Number of Vehicles}
The number of automobiles \( \text{numA} \) is generated using a Poisson random variable with parameter \( \lambda_A = 60 \).
The number of trucks \( \text{numT} \) is generated using a Poisson random variable with parameter \( \lambda_T = 12 \).

\subsection*{Weights Calculation}
The weight of each automobile is generated using a Gamma distribution with shape \( \alpha = 120 \) and scale \( \frac{1}{\lambda} = 10 \).

The weight of each truck is generated using a Gamma distribution with shape \( \alpha = 14 \) and scale \( \frac{1}{\lambda} = 1000 \).

\subsection*{Total Weight for Each Run}
For each Monte Carlo run, the total weight of all vehicles (automobiles and trucks) is calculated and stored in the \texttt{TotalWeight} vector.

\subsection*{Probability Estimation}
The probability of the bridge collapsing is estimated by counting the number of simulations where the total weight exceeds 250,000 kg (250 tons) and dividing by \( N \).

\[
p_{\text{est}} = \frac{\sum_{k=1}^{N} I(\text{TotalWeight}_k > 250000)}{N}
\]

\subsection*{Expected Value and Standard Deviation}
The expected total weight is the mean of the \texttt{TotalWeight} vector.

\[
\text{ExpectedWeight} = \frac{1}{N} \sum_{k=1}^{N} \text{TotalWeight}_k
\]

The standard deviation of the total weight is calculated using the standard deviation formula:

\[
\text{StdWeight} = \sqrt{\frac{1}{N-1} \sum_{k=1}^{N} (\text{TotalWeight}_k - \text{ExpectedWeight})^2}
\]

\section*{Output Example}
When we run the \texttt{hw4.m} script, the output will look something like this:

\begin{verbatim}
Estimated probability = 0.034217
Expected weight = 240123.487781
Standard deviation = 38504.637221
\end{verbatim}


\end{document}
